Per pair: student pages, pencils, a ruler, two nonstretch strings at least 40 cm long, and removable tape. Print at actual size. Main rectangles are \textbf{160 mm wide by 120 mm high}; verify the printed 100 mm calibration before using string or a ruler. The diagram in the key is deliberately smaller, not a measuring mat.

\textbf{Readiness.} K-1 needs right/up directions and side-by-side comparisons, with adult reading. Older entry requires addition of lengths; explaining the eight-turn lower bound needs control of a vertical-gap constraint. Children need not derive a square-root formula. Keep the three adult-led tables and offer the minimum-turn problem as exploration before proof.

\textbf{Launch, 4 minutes.} Let children bend string into a path, then follow the printed 40-by-30 mm example. Match A and B to the horizontal strip and C and D to the vertical strip. Show that string must follow every corner rather than cut across it. Ask, ``Can you make a very different route between the same two dots?''

\textbf{Flexible hour.} 0--5 launch; 5--22 build and compare; 22--25 walk an imagined right/up route; 25--43 narrow-strip or turn challenge; 43--55 compare explanations or allow left moves; 55--60 reset. Save a successful route before changing its rules.
